% ch16 — Apparent Horizon Finding
\chapter{Apparent Horizon Finding}
\label{ch:horizons}

Locating apparent horizons during a simulation provides real-time diagnostics:
black-hole masses, spins, and positions. This chapter develops the
expansion-scalar formulation for apparent horizons and the spectral
Newton-iteration algorithm that finds them.

\section{Trapped Surfaces and the Expansion Scalar}
\label{sec:ah-trapped}

We define marginally outer trapped surfaces (MOTS) via the vanishing of the
outgoing null expansion $\Theta_{(\ell)} = 0$ and explain their role as the
quasi-local boundary of a black hole.

\section{Spectral Parametrisation of the Horizon Surface}
\label{sec:ah-spectral}
\coderef{horizon\_finder}{rs}

We parametrise the horizon as a star-shaped surface $r = h(\theta,\phi)$
expanded in spherical harmonics, reducing the PDE to a finite set of spectral
coefficients.

\section{Newton Iteration}
\label{sec:ah-newton}

We linearise the expansion equation about a trial surface and solve the
resulting linear system by Newton-Raphson iteration, discussing convergence
criteria and the Jacobian structure.

\section{Horizon Properties: Mass, Spin, and Coordinate Centre}
\label{sec:ah-properties}

We compute the irreducible mass from the horizon area, extract the
dimensionless spin from the polar/equatorial circumferences, and track the
coordinate centroid for mesh-refinement box tracking.

\paragraph{Key References.}
The expansion-scalar formulation follows \cite{Thornburg2007}; the spectral
finder draws on \cite{Lin2007}; horizon diagnostics follow
\cite{Dreyer2003,Campanelli2006}.
